\documentclass[conference,onecolumn]{IEEEtran}
\usepackage{amsmath}
\usepackage{booktabs}
\usepackage{graphicx}
\usepackage{url}

% Draft checklist before submission:
% Replace student IDs, hardware details, and every TBD item with verified results.
% The existing results.csv contains one-epoch exploratory runs, NOT the final
% 30-epoch-budget study. Re-run all 21 experiments before completing the tables.
% Export the notebook figures and replace the boxed figure placeholders below.
% One-column working draft for reading and editing. For final submission,
% remove ",onecolumn" to restore the required IEEE conference layout;
% compile and check that the final version is <= 8 pages including references.

\newcommand{\tbd}{\textit{TBD}}

\title{How Architecture, Data Augmentation, and Optimiser Choice Affect CIFAR-10 Classification}
\author{\IEEEauthorblockN{STUDENT-ID-1 \quad STUDENT-ID-2}
\IEEEauthorblockA{PGR207 Deep Learning, Autumn 2026}}
% For a three-person group, add STUDENT-ID-3 to the author block and statement.

\begin{document}
\maketitle

\begin{abstract}
We investigate the effects of architecture, training-time data augmentation,
and optimiser choice on ten-class CIFAR-10 classification. A single shared
baseline anchors three one-factor-at-a-time studies: seven configurations
trained with three random seeds each. A fixed stratified training/validation
split, common evaluation protocol, and random weight initialisation make the
comparisons interpretable. We report test top-1 accuracy and macro-F1 as
mean $\pm$ standard deviation across seeds, alongside class-level errors,
training curves, model size, and training cost. \tbd{}: insert the principal
quantitative findings and a cautious conclusion after the full experiments.
\end{abstract}

\begin{IEEEkeywords}
CIFAR-10, image classification, convolutional neural networks, data
augmentation, optimisers, controlled experiments
\end{IEEEkeywords}

\section{Introduction}
CIFAR-10 contains small natural images assigned to ten mutually exclusive
classes~\cite{cifar10}. Architecture, augmentation, and optimisation can each
change classification performance, but changing several at once prevents us
from attributing an observed difference to a particular choice. We therefore
compare three levels of each factor against one shared baseline, while
holding data splits and the training and evaluation rules fixed.

We ask: (RQ1) how much does architecture change accuracy and macro-F1,
relative to model size and training cost? (RQ2) does stronger augmentation
improve generalisation under the same training budget? (RQ3) does the
optimiser change final performance or convergence speed? Our hypotheses,
specified before inspecting full-budget results, are that (H1) the
residual/dense models may outperform the plain CNN at greater cost; (H2)
crop/flip may improve generalisation, while stronger augmentation may need
more training; and (H3) optimisation will affect convergence, though a
shared learning rate could favour one optimiser. Differences must be
interpreted against the spread over three seeds, not a single best run.

\section{Related Work}
Residual connections facilitate training deeper image classifiers~\cite{resnet},
whereas DenseNet reuses features through dense connections~\cite{densenet}.
Together with a plain two-block CNN, these represent three distinct
architectural choices, rather than three widths of the same model.
Crop/flip is a standard image-classification augmentation; stronger policies
such as RandAugment broaden the set of image transformations~\cite{randaugment}.
A survey of augmentation methods motivates comparing progressively stronger
image-only transformations~\cite{augmentation_survey}. SGD with momentum,
Adam, and AdamW provide contrasting update rules~\cite{momentum,adam,adamw}.
AdamW decouples weight decay from its gradient-based update. All models here
start with random weights: pretrained ImageNet representations would obscure
differences induced by the choices under study. Our implementations use
PyTorch and TorchVision~\cite{pytorch,torchvision}.

\section{Methods}
\subsection{Data, preprocessing, and split}
We use the official CIFAR-10 split of 50,000 training and 10,000 test images
at $32\times32$ RGB resolution~\cite{cifar10}. A single class-stratified
split of the official training partition (random state 42) yields 45,000
training and 5,000 validation images; this split is reused across all runs.
Each of the ten classes contributes 4,500 training and 500 validation
images. Inputs are converted to floating point in $[0,1]$ and normalised
per channel using mean $(0.5,0.5,0.5)$ and standard deviation
$(0.5,0.5,0.5)$. Validation and test images receive only this deterministic
transform. We do not resize the images or use pretrained weights.

\subsection{Factor levels and controlled design}
The baseline (C1) is a plain CNN with no stochastic augmentation and SGD
with momentum. The CNN has two $3\times3$ convolution/ReLU blocks with
32 and 64 channels, each followed by $2\times2$ max pooling, then a
512-unit fully connected hidden layer and a ten-class output. For $32\times32$
inputs, ResNet-18 and DenseNet-121 replace their initial $7\times7$
stride-2 convolution with a $3\times3$ stride-1 convolution and remove
initial max pooling; both use ten output classes. The resulting model
parameter counts are 2,122,186 (CNN), 11,173,962 (ResNet-18), and
6,956,426 (DenseNet-121), counting all parameters in the notebook's
instantiated models.

The augmentation levels are (A0) no random transformation; (A1) random
crop to $32\times32$ with four-pixel padding, followed by random horizontal
flip; and (A2) A1 followed by RandAugment with two operations and magnitude
9~\cite{randaugment}. All levels share the same final normalisation, and
only training images are augmented. Optimiser levels are SGD with momentum
0.9, Adam, and AdamW. All use initial learning rate 0.01; SGD and Adam have
zero weight decay, while AdamW uses weight decay 0.01. Other optimiser
settings remain at PyTorch defaults~\cite{pytorch}.

\begin{table}[htbp]
\caption{The seven configurations. Each nonbaseline row changes exactly one factor relative to C1.}
\label{tab:design}
\centering
\normalsize
\setlength{\tabcolsep}{8pt}
\begin{tabular}{@{}llll@{}}
\toprule
ID & Architecture & Augmentation & Optimiser \\
\midrule
C1 & Plain CNN & A0 (none) & SGD + momentum \\
C2 & ResNet-18 & A0 & SGD + momentum \\
C3 & DenseNet-121 & A0 & SGD + momentum \\
C4 & Plain CNN & A1 (crop, flip) & SGD + momentum \\
C5 & Plain CNN & A2 (+RandAugment) & SGD + momentum \\
C6 & Plain CNN & A0 & Adam \\
C7 & Plain CNN & A0 & AdamW \\
\bottomrule
\end{tabular}
\end{table}

Table~\ref{tab:design} gives $1+2+2+2=7$ configurations. C1 is trained
once per seed and serves as the same reference in all three studies, for a
total of $7\times3=21$ runs. A row differing in two factors is deliberately
excluded. We use seeds 0, 1, and 2 to reinitialise model weights and
randomness for shuffling and augmentation. The fixed validation split is
not regenerated for each seed. Execution selects the Apple Metal
Performance Shaders (MPS) device when available, otherwise CPU.
\tbd{}: record the exact Mac chip, memory, PyTorch/TorchVision versions,
and whether every final run used MPS.

\subsection{Training and model selection}
All runs use batch size 64, cross-entropy loss, and the same policy of at
most 30 epochs. After each epoch, a ReduceLROnPlateau scheduler uses
validation loss, patience 3, and a learning-rate reduction factor of 0.1.
Training stops after six consecutive epochs without
improvement. The weights with lowest validation loss are restored before
a single test evaluation at the end of each run. Thus, actual epoch counts
may differ, but the stopping and scheduling criteria are identical. We log
sample-averaged training and validation loss and accuracy per epoch. Wall
time covers training and validation, excluding final test evaluation.

For the optimiser study we chose the assignment's
\emph{same-learning-rate} option: initial learning rate 0.01 for all three
optimisers. This strictly fixes the rate but is not necessarily equally
appropriate for SGD, Adam, and AdamW. We did not perform a per-optimiser
validation search; this limitation must be considered when interpreting
the optimiser comparison. Optimiser-specific weight decay is reported
above rather than silently treated as identical.

\subsection{Metrics and aggregation}
For each seed separately, test top-1 accuracy is the fraction of correctly
classified images, expressed as a percentage. Macro-F1 is the unweighted
mean of the ten class-wise F1 scores,
$\mathrm{macro\text{-}F1}=\frac{1}{10}\sum_{c=1}^{10}
\frac{2P_cR_c}{P_c+R_c}$, using zero for an undefined class score.
We calculate per-class precision $P_c$ and recall $R_c$ for the
configurations with lowest and highest \emph{mean best validation loss}
across seeds, showing the test confusion matrix of representative seed 0
for each. Selection of which examples to visualise does not change any
model or training decision. Summary tables give the mean $\pm$ sample
standard deviation ($n-1$ denominator) of the three independently
computed test metrics; predictions are never pooled across seeds.
The notebook writes one CSV row per run with metrics, parameter count,
training time, validation loss, and serialised epoch histories.

\section{Results}
% Replace all TBD entries using the full-budget 21-run CSV; reuse C1 unchanged.
Full-budget results are pending. The present one-epoch exploratory CSV is
not used for the comparisons below because it does not represent the stated
training budget or early-stopping policy. Tables~\ref{tab:arch}--\ref{tab:opt}
show the intended reporting format; all three must reuse identical C1 values.

\begin{table}[htbp]
\caption{Architecture study: test mean $\pm$ SD over seeds 0, 1, 2. Training seconds are the mean per run.}
\label{tab:arch}
\centering
\normalsize
\setlength{\tabcolsep}{8pt}
\begin{tabular}{@{}lrrrr@{}}
\toprule
ID/model & Params & Acc. (\%) & Macro-F1 & Time (s) \\
\midrule
C1 CNN & 2,122,186 & \tbd & \tbd & \tbd \\
C2 ResNet-18 & 11,173,962 & \tbd & \tbd & \tbd \\
C3 DenseNet-121 & 6,956,426 & \tbd & \tbd & \tbd \\
\bottomrule
\end{tabular}
\end{table}

\begin{table}[htbp]
\caption{Augmentation study: test mean $\pm$ SD over three seeds. C1 repeats Table~\ref{tab:arch}.}
\label{tab:aug}
\centering
\normalsize
\setlength{\tabcolsep}{8pt}
\begin{tabular}{@{}lrrr@{}}
\toprule
ID/strategy & Acc. (\%) & Macro-F1 & Time (s) \\
\midrule
C1 none & \tbd & \tbd & \tbd \\
C4 crop + flip & \tbd & \tbd & \tbd \\
C5 +RandAugment & \tbd & \tbd & \tbd \\
\bottomrule
\end{tabular}
\end{table}

\begin{table}[htbp]
\caption{Optimiser study: test mean $\pm$ SD over three seeds. C1 repeats Table~\ref{tab:arch}.}
\label{tab:opt}
\centering
\normalsize
\setlength{\tabcolsep}{8pt}
\begin{tabular}{@{}lrrr@{}}
\toprule
ID/optimiser & Acc. (\%) & Macro-F1 & Time (s) \\
\midrule
C1 SGD + momentum & \tbd & \tbd & \tbd \\
C6 Adam & \tbd & \tbd & \tbd \\
C7 AdamW & \tbd & \tbd & \tbd \\
\bottomrule
\end{tabular}
\end{table}

\tbd{}: identify the two configurations by mean best validation loss and
transcribe the seed-0 test precision/recall of all ten classes into
Table~\ref{tab:classes}. Report the two corresponding confusion matrices;
discuss specific off-diagonal pairs rather than only the diagonal.

\begin{table}[htbp]
\caption{Per-class test precision and recall for validation-selected best (B) and worst (W) configurations; representative seed 0. Replace the labels and values after the full runs.}
\label{tab:classes}
\centering
\normalsize
\setlength{\tabcolsep}{8pt}
\begin{tabular}{@{}lrrrr@{}}
\toprule
Class & $P_{\rm B}$ & $R_{\rm B}$ & $P_{\rm W}$ & $R_{\rm W}$ \\
\midrule
airplane   & \tbd & \tbd & \tbd & \tbd \\
automobile & \tbd & \tbd & \tbd & \tbd \\
bird       & \tbd & \tbd & \tbd & \tbd \\
cat        & \tbd & \tbd & \tbd & \tbd \\
deer       & \tbd & \tbd & \tbd & \tbd \\
dog        & \tbd & \tbd & \tbd & \tbd \\
frog       & \tbd & \tbd & \tbd & \tbd \\
horse      & \tbd & \tbd & \tbd & \tbd \\
ship       & \tbd & \tbd & \tbd & \tbd \\
truck      & \tbd & \tbd & \tbd & \tbd \\
\bottomrule
\end{tabular}
\end{table}

\begin{figure}[htbp]
\centering
\fbox{\parbox[c][0.65in][c]{0.88\columnwidth}{\centering TBD: insert best and worst seed-0 confusion matrices}}
\caption{Test confusion matrices for the best and worst configurations,
selected by mean best validation loss. Specify both configuration IDs.}
\label{fig:confusion}
\end{figure}

\begin{figure}[htbp]
\centering
\fbox{\parbox[c][0.65in][c]{0.88\columnwidth}{\centering TBD: insert training/validation curves}}
\caption{Representative seed-0 training and validation loss/accuracy;
include optimiser convergence curves when discussing RQ3.}
\label{fig:curves}
\end{figure}

\section{Discussion and Conclusion}
\tbd{}: for RQ1--RQ3, insert differences against C1 in percentage points,
the corresponding across-seed SD, and an interpretation of whether each
difference exceeds the observed noise floor. Rank the three factors by
effect size, state whether H1--H3 were supported, and compare accuracy
per parameter and per unit training time where useful. Explain the
classes most often confused in Fig.~\ref{fig:confusion} and use
Fig.~\ref{fig:curves} to distinguish convergence from final performance.
If a difference is comparable to seed variation, report an inconclusive
comparison rather than a winner.

This design measures each factor at only one setting of the other two;
it cannot establish interactions. The fixed training budget, three seeds,
small $32\times32$ images, and many comparisons against one test set limit
generalisation. The common learning rate may bias the optimiser comparison,
and training/validation wall time also includes data augmentation and
hardware overhead. \tbd{}: replace this final sentence with the principal
evidence-based conclusion once the results are complete.

\section*{Contribution Statement}
STUDENT-ID-1: \tbd{} (training, analysis, and writing contributions).
STUDENT-ID-2: \tbd{} (training, analysis, and writing contributions).
% Add a third student ID and contribution sentence if applicable.
All members reviewed and can explain the full experimental pipeline and
manuscript.

\section*{AI Usage Statement}
OpenCode (an AI assistant) was used to assist with notebook coding,
debugging, and preparation of an initial LaTeX manuscript draft. The
authors are responsible for checking the implementation, citations,
reported numbers, and final interpretations. \tbd{}: add any other AI
tools actually used and confirm this statement with the group.

\begin{thebibliography}{00}
\bibitem{cifar10} A. Krizhevsky, ``Learning multiple layers of features from tiny images,'' University of Toronto, Tech. Rep., 2009. [Online]. Available: \url{https://www.cs.toronto.edu/~kriz/cifar.html}
\bibitem{torchvision} TorchVision maintainers and contributors, ``TorchVision: PyTorch's computer vision library,'' 2016. [Online]. Available: \url{https://github.com/pytorch/vision}
\bibitem{resnet} K. He, X. Zhang, S. Ren, and J. Sun, ``Deep residual learning for image recognition,'' in \emph{Proc. IEEE Conf. Computer Vision and Pattern Recognition}, 2016, pp. 770--778, doi: 10.1109/CVPR.2016.90.
\bibitem{densenet} G. Huang, Z. Liu, L. van der Maaten, and K. Q. Weinberger, ``Densely connected convolutional networks,'' in \emph{Proc. IEEE Conf. Computer Vision and Pattern Recognition}, 2017, pp. 4700--4708, doi: 10.1109/CVPR.2017.243.
\bibitem{randaugment} E. D. Cubuk, B. Zoph, J. Shlens, and Q. V. Le, ``RandAugment: Practical automated data augmentation with a reduced search space,'' in \emph{Proc. IEEE/CVF Conf. Computer Vision and Pattern Recognition Workshops}, 2020, pp. 702--703, doi: 10.1109/CVPRW50498.2020.00359.
\bibitem{augmentation_survey} C. Shorten and T. M. Khoshgoftaar, ``A survey on image data augmentation for deep learning,'' \emph{Journal of Big Data}, vol. 6, no. 60, 2019, doi: 10.1186/s40537-019-0197-0.
\bibitem{momentum} I. Sutskever, J. Martens, G. Dahl, and G. Hinton, ``On the importance of initialization and momentum in deep learning,'' in \emph{Proc. International Conf. Machine Learning}, 2013, pp. 1139--1147.
\bibitem{adam} D. P. Kingma and J. Ba, ``Adam: A method for stochastic optimization,'' in \emph{Proc. International Conf. Learning Representations}, 2015.
\bibitem{adamw} I. Loshchilov and F. Hutter, ``Decoupled weight decay regularization,'' in \emph{Proc. International Conf. Learning Representations}, 2019.
\bibitem{pytorch} A. Paszke \emph{et al.}, ``PyTorch: An imperative style, high-performance deep learning library,'' in \emph{Advances in Neural Information Processing Systems 32}, 2019, pp. 8024--8035.
\end{thebibliography}

\end{document}
